\clearpage
\VerdierPhysicalPageAnchor{0130}%
\begingroup
\small
\setlength{\emergencystretch}{1.5em}%
\setlength{\abovedisplayskip}{2pt plus 1pt minus 1pt}%
\setlength{\belowdisplayskip}{2pt plus 1pt minus 1pt}%

\VerdierSemanticAnchor{ch2-2-1-3-2-sm2-completion}
\noindent se complète en un diagramme commutatif :

\VerdierSemanticAnchor{ch2-2-1-3-2-completion-diagram}
\begin{center}
  $t\in S$\qquad
  \begin{tikzcd}[column sep=4.0em,row sep=2.5em,baseline=(current bounding box.center)]
    Z' \arrow[r,dashed,"g"] \arrow[d,dashed,"t"'] &
    X' \arrow[d,"s"] \\
    Z \arrow[r,"f"'] & X
  \end{tikzcd}
  \quad ,
\end{center}

\VerdierSemanticAnchor{ch2-2-1-3-2-opposite-reduction}
\noindent car alors la deuxième assertion de SM2) s'obtiendra en passant à
la catégorie opposée.

\VerdierSemanticAnchor{ch2-2-1-3-2-triangle-construction}
Par TRI) et TRII), on sait qu'il existe un triangle distingué
$(X,Y,Z[1],u,v,f[1])$ et, par TRI), on sait qu'il existe un triangle distingué
$(X',Y,Z',us,v',w')$. On a alors un diagramme commutatif :

\VerdierSemanticAnchor{ch2-2-1-3-2-us-square}
\begin{center}
  \begin{tikzcd}[column sep=4.2em,row sep=2.5em]
    X' \arrow[r,"us"] \arrow[d,"s"'] &
    Y \arrow[d,"\operatorname{id}_Y"] \\
    X \arrow[r,"u"'] & Y
  \end{tikzcd}
  \quad .
\end{center}

\VerdierSemanticAnchor{ch2-2-1-3-2-sm6-application}
\noindent Le morphisme $\operatorname{id}_Y$ appartient à $S$ (SM1) et, par
SM6), on obtient un morphisme $s'\colon Z'\to Z[1]$ appartenant à $S$ tel que
$(s,\operatorname{id}_Y,s')$ soit un morphisme de triangles. En particulier,
le diagramme ci-après est commutatif :

\VerdierSemanticAnchor{ch2-2-1-3-2-shift-square}
\begin{center}
  \begin{tikzcd}[column sep=4.2em,row sep=2.5em]
    Z' \arrow[r,"{w'}"] \arrow[d,"{s'}"'] &
    X'[1] \arrow[d,"{s[1]}"] \\
    Z[1] \arrow[r,"{f[1]}"'] & X[1]
  \end{tikzcd}
  \quad ,
\end{center}

\VerdierSemanticAnchor{ch2-2-1-3-2-proof-complete}
\noindent d'où, en utilisant SM5), la première assertion de SM2).

\VerdierSemanticAnchor{ch2-2-1-4}
\medskip
\noindent\textbf{2.1.4.}\quad
\VerdierSemanticAnchor{ch2-2-1-4-convention}
\emph{Les systèmes multiplicatifs des catégories triangulées seront, sauf
mention expresse du contraire, compatibles avec la triangulation. Ils seront
donc appelés simplement et par abus de langage, systèmes multiplicatifs.}

\VerdierSemanticAnchor{ch2-2-1-5}
\medskip
\noindent\textbf{2.1.5.}\quad
\VerdierSemanticAnchor{ch2-2-1-5-full-triangulated-subcategory}
Soit $\mathcal D$ une catégorie triangulée. Une \emph{sous-catégorie
triangulée pleine} de $\mathcal D$ est une sous-catégorie pleine de
$\mathcal D$ stable par translations, munie d'une structure triangulée telle
que le foncteur d'inclusion soit un foncteur exact (pour la commutation stricte
aux translations).

\VerdierSemanticAnchor{ch2-2-1-5-additive-subcategory}
Une sous-catégorie triangulée pleine de $\mathcal D$ est une sous-catégorie
additive de $\mathcal D$ (1.2.7).

\VerdierSemanticAnchor{ch2-2-1-5-boundary}
Une sous-catégorie
\endgroup
